% Section prête à insérer dans rapport.tex, juste avant « \section{Limites} », avec : % Section prête à insérer dans rapport.tex, juste avant « \section{Limites} », avec : % Section prête à insérer dans rapport.tex, juste avant « \section{Limites} », avec : % Section prête à insérer dans rapport.tex, juste avant « \section{Limites} », avec : \input{section_apres_coup}
% Elle n'est pas encore incluse : le rapport n'a pas pu être recompilé sur cette machine (pas de LaTeX installé).
% Tous les chiffres viennent de runs/exploration_*.json (python -m src.banc_essai, src.banc_essai_jour,
% src.banc_essai_long, src.volatilite).

\section{Après le test : vérifications et second volet}

Tout ce qui suit a été fait \textbf{après} l'ouverture du test. Ces essais éclairent le résultat de la section
précédente ; ils ne le remplacent pas, et les fichiers figés n'ont pas été modifiés.

\paragraph{Un autre modèle aurait-il fait mieux ?} Nous avons réentraîné chaque mois vingt modèles (logistiques,
forêts, LightGBM, réseau de neurones, AR(1), lissage exponentiel\ldots), sur tout le passé ou sur les 60 derniers mois.
Sur les 33 mois du test, 3 configurations sur 38 passent sous la log-loss de la constante, avec une avance
($-0{,}009 \pm 0{,}039$ pour la meilleure) quatre fois plus petite que son erreur type. Le modèle le mieux classé sur
2016-2023 (15 plus proches voisins) finit parmi les moins bons sur le test. Sur la base journalière (4\,118 séances),
l'AUC du sens du lendemain vaut 0,49 sur 668 séances de test. Faire commencer l'apprentissage en 2000 plutôt qu'en 2010
dégrade la constante (0,667 contre 0,660). Le blocage vient donc de l'information disponible, pas du choix du modèle.

\paragraph{Un piège d'horaires.} Dans la base journalière, la ligne d'un jour contient la clôture de Wall Street du
même jour, connue après celle de Tokyo. L'utiliser telle quelle donne une AUC de 0,71, sans décision possible derrière.
Toutes les séries étrangères ont donc été décalées d'un jour.

\paragraph{Second volet : prévoir l'agitation.} Avec les mêmes mois, le même découpage et un réentraînement mensuel,
nous avons changé de cible : la volatilité réalisée du mois suivant, et la classe « mois agité » (volatilité au-dessus
de 20~\%).

\begin{table}[H]\centering\small
\begin{tabular}{lcc}\toprule
 & 2016-2023 (95 mois) & Test 2024-2026 (33 mois) \\\midrule
Volatilité : $R^2$ hors échantillon, régression (vol. 1, 3, 6 mois, VIX, rendement) & 0,23 & 0,18 \\
Volatilité : $R^2$ hors échantillon, GARCH(1,1) & 0,03 & 0,32 \\
Mois agité : AUC de la régression logistique & 0,73 & 0,71 \\
Mois agité : écart de log-loss avec la constante & $-0{,}074 \pm 0{,}041$ & $-0{,}065 \pm 0{,}088$ \\\bottomrule
\end{tabular}\end{table}

Le classement des mois est stable d'une période à l'autre (AUC proche de 0,7, contre 0,5 pour le sens du marché),
mais sur 33 mois l'avantage en log-loss reste inférieur à son erreur type, et le classement des méthodes de
prévision change selon la période. Une position dosée selon la volatilité prévue (part investie $= 15~\%/\hat\sigma$,
plafonnée à 100~\%) ramène le pire mois de $-13{,}2~\%$ à $-9{,}7~\%$ et le rendement de $+82{,}3~\%$ à $+67{,}7~\%$ :
rapporté au risque, le résultat est le même (1,00 contre 0,98). Ce volet sert à annoncer le risque ; il ne fournit pas
une stratégie meilleure que « toujours acheter, couvert ».

% Elle n'est pas encore incluse : le rapport n'a pas pu être recompilé sur cette machine (pas de LaTeX installé).
% Tous les chiffres viennent de runs/exploration_*.json (python -m src.banc_essai, src.banc_essai_jour,
% src.banc_essai_long, src.volatilite).

\section{Après le test : vérifications et second volet}

Tout ce qui suit a été fait \textbf{après} l'ouverture du test. Ces essais éclairent le résultat de la section
précédente ; ils ne le remplacent pas, et les fichiers figés n'ont pas été modifiés.

\paragraph{Un autre modèle aurait-il fait mieux ?} Nous avons réentraîné chaque mois vingt modèles (logistiques,
forêts, LightGBM, réseau de neurones, AR(1), lissage exponentiel\ldots), sur tout le passé ou sur les 60 derniers mois.
Sur les 33 mois du test, 3 configurations sur 38 passent sous la log-loss de la constante, avec une avance
($-0{,}009 \pm 0{,}039$ pour la meilleure) quatre fois plus petite que son erreur type. Le modèle le mieux classé sur
2016-2023 (15 plus proches voisins) finit parmi les moins bons sur le test. Sur la base journalière (4\,118 séances),
l'AUC du sens du lendemain vaut 0,49 sur 668 séances de test. Faire commencer l'apprentissage en 2000 plutôt qu'en 2010
dégrade la constante (0,667 contre 0,660). Le blocage vient donc de l'information disponible, pas du choix du modèle.

\paragraph{Un piège d'horaires.} Dans la base journalière, la ligne d'un jour contient la clôture de Wall Street du
même jour, connue après celle de Tokyo. L'utiliser telle quelle donne une AUC de 0,71, sans décision possible derrière.
Toutes les séries étrangères ont donc été décalées d'un jour.

\paragraph{Second volet : prévoir l'agitation.} Avec les mêmes mois, le même découpage et un réentraînement mensuel,
nous avons changé de cible : la volatilité réalisée du mois suivant, et la classe « mois agité » (volatilité au-dessus
de 20~\%).

\begin{table}[H]\centering\small
\begin{tabular}{lcc}\toprule
 & 2016-2023 (95 mois) & Test 2024-2026 (33 mois) \\\midrule
Volatilité : $R^2$ hors échantillon, régression (vol. 1, 3, 6 mois, VIX, rendement) & 0,23 & 0,18 \\
Volatilité : $R^2$ hors échantillon, GARCH(1,1) & 0,03 & 0,32 \\
Mois agité : AUC de la régression logistique & 0,73 & 0,71 \\
Mois agité : écart de log-loss avec la constante & $-0{,}074 \pm 0{,}041$ & $-0{,}065 \pm 0{,}088$ \\\bottomrule
\end{tabular}\end{table}

Le classement des mois est stable d'une période à l'autre (AUC proche de 0,7, contre 0,5 pour le sens du marché),
mais sur 33 mois l'avantage en log-loss reste inférieur à son erreur type, et le classement des méthodes de
prévision change selon la période. Une position dosée selon la volatilité prévue (part investie $= 15~\%/\hat\sigma$,
plafonnée à 100~\%) ramène le pire mois de $-13{,}2~\%$ à $-9{,}7~\%$ et le rendement de $+82{,}3~\%$ à $+67{,}7~\%$ :
rapporté au risque, le résultat est le même (1,00 contre 0,98). Ce volet sert à annoncer le risque ; il ne fournit pas
une stratégie meilleure que « toujours acheter, couvert ».

% Elle n'est pas encore incluse : le rapport n'a pas pu être recompilé sur cette machine (pas de LaTeX installé).
% Tous les chiffres viennent de runs/exploration_*.json (python -m src.banc_essai, src.banc_essai_jour,
% src.banc_essai_long, src.volatilite).

\section{Après le test : vérifications et second volet}

Tout ce qui suit a été fait \textbf{après} l'ouverture du test. Ces essais éclairent le résultat de la section
précédente ; ils ne le remplacent pas, et les fichiers figés n'ont pas été modifiés.

\paragraph{Un autre modèle aurait-il fait mieux ?} Nous avons réentraîné chaque mois vingt modèles (logistiques,
forêts, LightGBM, réseau de neurones, AR(1), lissage exponentiel\ldots), sur tout le passé ou sur les 60 derniers mois.
Sur les 33 mois du test, 3 configurations sur 38 passent sous la log-loss de la constante, avec une avance
($-0{,}009 \pm 0{,}039$ pour la meilleure) quatre fois plus petite que son erreur type. Le modèle le mieux classé sur
2016-2023 (15 plus proches voisins) finit parmi les moins bons sur le test. Sur la base journalière (4\,118 séances),
l'AUC du sens du lendemain vaut 0,49 sur 668 séances de test. Faire commencer l'apprentissage en 2000 plutôt qu'en 2010
dégrade la constante (0,667 contre 0,660). Le blocage vient donc de l'information disponible, pas du choix du modèle.

\paragraph{Un piège d'horaires.} Dans la base journalière, la ligne d'un jour contient la clôture de Wall Street du
même jour, connue après celle de Tokyo. L'utiliser telle quelle donne une AUC de 0,71, sans décision possible derrière.
Toutes les séries étrangères ont donc été décalées d'un jour.

\paragraph{Second volet : prévoir l'agitation.} Avec les mêmes mois, le même découpage et un réentraînement mensuel,
nous avons changé de cible : la volatilité réalisée du mois suivant, et la classe « mois agité » (volatilité au-dessus
de 20~\%).

\begin{table}[H]\centering\small
\begin{tabular}{lcc}\toprule
 & 2016-2023 (95 mois) & Test 2024-2026 (33 mois) \\\midrule
Volatilité : $R^2$ hors échantillon, régression (vol. 1, 3, 6 mois, VIX, rendement) & 0,23 & 0,18 \\
Volatilité : $R^2$ hors échantillon, GARCH(1,1) & 0,03 & 0,32 \\
Mois agité : AUC de la régression logistique & 0,73 & 0,71 \\
Mois agité : écart de log-loss avec la constante & $-0{,}074 \pm 0{,}041$ & $-0{,}065 \pm 0{,}088$ \\\bottomrule
\end{tabular}\end{table}

Le classement des mois est stable d'une période à l'autre (AUC proche de 0,7, contre 0,5 pour le sens du marché),
mais sur 33 mois l'avantage en log-loss reste inférieur à son erreur type, et le classement des méthodes de
prévision change selon la période. Une position dosée selon la volatilité prévue (part investie $= 15~\%/\hat\sigma$,
plafonnée à 100~\%) ramène le pire mois de $-13{,}2~\%$ à $-9{,}7~\%$ et le rendement de $+82{,}3~\%$ à $+67{,}7~\%$ :
rapporté au risque, le résultat est le même (1,00 contre 0,98). Ce volet sert à annoncer le risque ; il ne fournit pas
une stratégie meilleure que « toujours acheter, couvert ».

% Elle n'est pas encore incluse : le rapport n'a pas pu être recompilé sur cette machine (pas de LaTeX installé).
% Tous les chiffres viennent de runs/exploration_*.json (python -m src.banc_essai, src.banc_essai_jour,
% src.banc_essai_long, src.volatilite).

\section{Après le test : vérifications et second volet}

Tout ce qui suit a été fait \textbf{après} l'ouverture du test. Ces essais éclairent le résultat de la section
précédente ; ils ne le remplacent pas, et les fichiers figés n'ont pas été modifiés.

\paragraph{Un autre modèle aurait-il fait mieux ?} Nous avons réentraîné chaque mois vingt modèles (logistiques,
forêts, LightGBM, réseau de neurones, AR(1), lissage exponentiel\ldots), sur tout le passé ou sur les 60 derniers mois.
Sur les 33 mois du test, 3 configurations sur 38 passent sous la log-loss de la constante, avec une avance
($-0{,}009 \pm 0{,}039$ pour la meilleure) quatre fois plus petite que son erreur type. Le modèle le mieux classé sur
2016-2023 (15 plus proches voisins) finit parmi les moins bons sur le test. Sur la base journalière (4\,118 séances),
l'AUC du sens du lendemain vaut 0,49 sur 668 séances de test. Faire commencer l'apprentissage en 2000 plutôt qu'en 2010
dégrade la constante (0,667 contre 0,660). Le blocage vient donc de l'information disponible, pas du choix du modèle.

\paragraph{Un piège d'horaires.} Dans la base journalière, la ligne d'un jour contient la clôture de Wall Street du
même jour, connue après celle de Tokyo. L'utiliser telle quelle donne une AUC de 0,71, sans décision possible derrière.
Toutes les séries étrangères ont donc été décalées d'un jour.

\paragraph{Second volet : prévoir l'agitation.} Avec les mêmes mois, le même découpage et un réentraînement mensuel,
nous avons changé de cible : la volatilité réalisée du mois suivant, et la classe « mois agité » (volatilité au-dessus
de 20~\%).

\begin{table}[H]\centering\small
\begin{tabular}{lcc}\toprule
 & 2016-2023 (95 mois) & Test 2024-2026 (33 mois) \\\midrule
Volatilité : $R^2$ hors échantillon, régression (vol. 1, 3, 6 mois, VIX, rendement) & 0,23 & 0,18 \\
Volatilité : $R^2$ hors échantillon, GARCH(1,1) & 0,03 & 0,32 \\
Mois agité : AUC de la régression logistique & 0,73 & 0,71 \\
Mois agité : écart de log-loss avec la constante & $-0{,}074 \pm 0{,}041$ & $-0{,}065 \pm 0{,}088$ \\\bottomrule
\end{tabular}\end{table}

Le classement des mois est stable d'une période à l'autre (AUC proche de 0,7, contre 0,5 pour le sens du marché),
mais sur 33 mois l'avantage en log-loss reste inférieur à son erreur type, et le classement des méthodes de
prévision change selon la période. Une position dosée selon la volatilité prévue (part investie $= 15~\%/\hat\sigma$,
plafonnée à 100~\%) ramène le pire mois de $-13{,}2~\%$ à $-9{,}7~\%$ et le rendement de $+82{,}3~\%$ à $+67{,}7~\%$ :
rapporté au risque, le résultat est le même (1,00 contre 0,98). Ce volet sert à annoncer le risque ; il ne fournit pas
une stratégie meilleure que « toujours acheter, couvert ».
